\subsection{连接同态与同调正合列}

在本节中，我们考虑复形的一个正合列

\[
0 \longrightarrow \mathrm{C} ^ {\prime} \stackrel {u} {\longrightarrow} \mathrm{C} \stackrel {v} {\longrightarrow} \mathrm{C} ^ {\prime \prime} \longrightarrow 0;
\]

用同一字母d表示C、C'和C''的微分\,'.

设 \(\Gamma\) 为由 \(x \in \mathbf{C}\) 中满足 \(dx \in \operatorname{Im}(u)\) 的元素组成的集合；对于 \(x \in \Gamma\) ，我们有

\[
d \left(^ {- 1} u ^ {1} (d x)\right) = ^ {- 1} u ^ {1} (d d (x)) = 0,
\]

因此 \(\overline{u}^{1}(dx)\in\mathbf{Z}(\mathbf{C}^{\prime})\) ；我们还有 \(dv(x)=v(dx)\in\operatorname{Im}(v\circ u)=0\) ，因此 \(v(x)\in\mathbf{Z}(\mathbf{C}^{\prime\prime})\) ；于是考虑线性映射 \(\varphi:\Gamma\to\mathrm{H}(\mathbf{C}^{\prime\prime})\times\mathrm{H}(\mathbf{C}^{\prime})\) ，它将每个元素 \(x\in\Gamma\) 映到 \((v(x),\overline{u}^{1}(dx))\) 的类 .

引理1. —— \(\varphi(\Gamma)\) （ \(\Gamma\) 在 \(\mathrm{H}(\mathbf{C}^{\prime\prime}) \times \mathrm{H}(\mathbf{C}^{\prime})\) 中的像）是从 \(\mathrm{H}(\mathbf{C}^{\prime\prime})\) 到 \(\mathrm{H}(\mathbf{C}^{\prime})\) 的一个次数为 -1 的分次 A-同态的图 .

\begin{enumerate}
\def\labelenumi{\alph{enumi})}
\item
  若 \(x \in \Gamma\) 且 \(v(x) \in \mathbf{B}(\mathbf{C}'')\) ，则 \(\bar{u}^1(dx) \in \mathbf{B}(\mathbf{C}')\) ：确实存在 \(z'' \in \mathbf{C}''\) 使得 \(v(x) = dz''\) ，则 \(z \in \mathbf{C}\) 使得 \(z'' = v(z)\) ，因此 \(v(x) = v(dz)\) ，则 \(t' \in \mathbf{C}'\) 使得 \(x - dz = u(t')\) ，它给出 \(dx = u(dt')\) ，因此 \(\bar{u}^1(dx) = dt' \in \mathbf{B}(\mathbf{C}')\) .
\item
  \(\mathbf{Z}(\mathbf{C}^{\prime \prime})\) 中的每个元素都是 \(v\) 作用于某个元素 \(x\) （属于 \(\mathbf{C}\) 且满足 \(v(dx) = 0\) ，即\(dx \in \operatorname{Im} u\) ，亦即 \(x \in \Gamma\) ）的像 .
\item
  由a)和b)可知 \(\varphi(\Gamma)\) 确实是一个函数图；由于 \(\varphi\) 是双齐次的，且双次数为 \((0, -1)\) ，这完成了证明。
\end{enumerate}

由此定义的从 H(C'') 到 H(C') 的次数为 -1 的分次 A-同态称为相对于正合列 (u, v) 的连接同态；它记为 \(\partial(u, v)\) 或 \(\partial_{u,v}\) 或简记为 \(\partial\) 。其齐次分量记为

\[
\partial_ {n} (u, v): \mathrm{H} _ {n} (\mathbf {C} ^ {\prime \prime}) \rightarrow \mathrm{H} _ {n - 1} (\mathbf {C} ^ {\prime}) \quad \text { and } \quad \partial^ {n} (u, v): \mathrm{H} ^ {n} (\mathbf {C} ^ {\prime \prime}) \rightarrow \mathrm{H} ^ {n + 1} (\mathbf {C} ^ {\prime}).
\]

根据定义，为了构造类 \(\alpha\in\mathrm{H}_{n}(\mathbf{C}^{\prime\prime})\) 在 \(\partial\) 下的像，我们选取一个闭链 \(z^{\prime\prime}\in\mathbf{Z}_{n}(\mathbf{C}^{\prime\prime})\) （它代表类 \(\alpha\) ），然后选取一个元素 \(x\) （属于 \(\mathbf{C}_{n}\) ），使得 \(v(x)=z^{\prime\prime}\) ；于是 \(dx\) 具有 \(u(t')\) 的形式， \(t^{\prime}\in\mathbf{C}_{n-1}^{\prime}\) ，而 \(\partial(\alpha)\) 是 \(t'\) 的同调类 .

用对应的话来说， \(\partial_{n}(u,v)\) 因此是由对应 \(u_{n-1}^{-1}\circ d_{n}\circ v_{n}^{-1}\) （介于 \(\mathbf{C}_{n}^{\prime\prime}\) 与 \(\mathbf{C}_{n-1}^{\prime}\) 之间）通过先过渡到子集 \(\mathbf{Z}_{n}(\mathbf{C}^{\prime\prime})\) 与 \(\mathbf{Z}_{n-1}(\mathbf{C}^{\prime})\) 、再过渡到它们的商 \(\mathbf{H}_{n}(\mathbf{C}^{\prime\prime})\) 与 \(\mathbf{H}_{n-1}(\mathbf{C}^{\prime})\) 而得到的。这特别表明，如果把 \(\mathbf{C},\mathbf{C}^{\prime},\mathbf{C}^{\prime\prime},u,v\) 替换为 \(\mathbf{C}(p),\mathbf{C}^{\prime}(p),\mathbf{C}^{\prime\prime}(p),u(p),v(p)\) ，则有

\[
\partial (u (p), v (p)) = (- 1) ^ {p} \partial (u, v);\tag{4}
\]

同样地，如果 \(\lambda\) 与 \(\mu\) 是 A 的中心的两个可逆元素，则有

\[
\partial (\lambda u, \mu v) = \lambda^ {- 1} \mu^ {- 1} \partial (u, v).\tag{5}
\]

还可以将 \(\partial(u,v)\) 与蛇形图（X，第 4 页）联系起来。根据 loc. cit.，命题 2，序列

\[
0 \longrightarrow Z _ {n} (C ^ {\prime}) \xrightarrow {Z _ {n} (u)} Z _ {n} (C) \xrightarrow {Z _ {n} (v)} Z _ {n} (C ^ {\prime \prime})
\]

以及

\[
\mathrm{C} _ {n} ^ {\prime} / \mathrm{B} _ {n} (\mathrm{C} ^ {\prime}) \xrightarrow {\overline {{{u}}} _ {n}} \mathrm{C} _ {n} / \mathrm{B} _ {n} (\mathrm{C}) \xrightarrow {\overline {{{v}}} _ {n}} \mathrm{C} _ {n} ^ {\prime \prime} / \mathrm{B} _ {n} (\mathrm{C} ^ {\prime}) \longrightarrow 0,
\]

其中 \(\overline{u}_{n}\) 与 \(\overline{v}_{n}\) 分别由 \(u\) 与 \(v\) 诱导，都是正合的。利用典范正合列 \((\mathrm{V}_{n})\) ，可以得到一个具有正合行和正合列的交换图

\[
\begin{array}{c c c c} 0 & 0 & 0 \\ \downarrow & \mathrm{H} _ {n} (\mathrm{C} ^ {\prime}) & \xrightarrow {\mathrm{H} _ {n} (u)} & \mathrm{H} _ {n} (\mathrm{C}) \\ \downarrow & \mathrm{C} _ {n} ^ {\prime} / \mathrm{B} _ {n} (\mathrm{C} ^ {\prime}) & \xrightarrow {\bar {u} _ {n}} & \mathrm{C} _ {n} / \mathrm{B} _ {n} (\mathrm{C}) \\ \downarrow & 0 \longrightarrow Z _ {n - 1} (\mathrm{C} ^ {\prime}) & \xrightarrow {Z _ {n - 1} (u)} & Z _ {n - 1} (\mathrm{C}) \\ \downarrow & H _ {n - 1} (\mathrm{C} ^ {\prime}) & \xrightarrow {\mathrm{H} _ {n - 1} (u)} & H _ {n - 1} (\mathrm{C}) \\ \downarrow & 0 & 0 & 0 \end{array}
\]

与这个图相关联的同态 \(\mathrm{H}_{n}(\mathbf{C}^{\prime \prime})\to \mathrm{H}_{n - 1}(\mathbf{C}^{\prime})\) （loc. cit.，命题 2，(iii)）按构造与 \(\partial_n(u,v)\) 一致。这也意味着同态序列 \((\mathrm{H}_n(u),\mathrm{H}_n(v),\partial_n(u,v),\mathrm{H}_{n - 1}(u),\mathrm{H}_{n - 1}(v))\) 是正合的；因此：

定理1. --- A-模同态的无穷序列

\[
\begin{array}{r l} \dots & \to \mathrm{H} _ {n + 1} (\mathrm{C} ^ {\prime \prime}) \xrightarrow {\partial_ {n + 1} (u , v)} \mathrm{H} _ {n} (\mathrm{C} ^ {\prime}) \xrightarrow {\mathrm{H} _ {n} (u)} \mathrm{H} _ {n} (\mathrm{C}) \xrightarrow {\mathrm{H} _ {n} (v)} \mathrm{H} _ {n} (\mathrm{C} ^ {\prime \prime}) \\ & \xrightarrow {\partial_ {n} (u , v)} \mathrm{H} _ {n - 1} (\mathrm{C} ^ {\prime}) \xrightarrow {\mathrm{H} _ {n - 1} (u)} \mathrm{H} _ {n - 1} (\mathrm{C}) \xrightarrow {\mathrm{H} _ {n - 1} (v)} \mathrm{H} _ {n - 1} (\mathrm{C} ^ {\prime \prime}) \xrightarrow {\partial_ {n - 1} (u , v)} \mathrm{H} _ {n - 2} (\mathrm{C} ^ {\prime}) \to \dots \end{array}
\]

是正合的。

这个序列称为与正合列 \((u, v)\) 相关联的同调正合列；它有时写成 A-模的正合三角形的形式

\[
\begin{array}{c} \mathrm{H(C)} \\ \mathrm{H(u)} \nearrow \searrow \mathrm{H(v)} \\ \mathrm {H(C^ {\prime})} \xleftarrow [ \partial (u , v) ]{} \mathrm {H(C^{\prime\prime})}. \end{array}
\]

推论1. --- 若复形 C、C'、C'' 中有两个具有零同调，则第三个也具有零同调。要使 u（相应地 v）为拟同构，必要且充分条件是 C''（相应地 C'）具有零同调。要使 \(\partial(u, v)\) 是双射的，必要且充分条件是 C 具有零同调。

推论2. —— 设u是一个复形态射。如果Ker u和Coker u具有零同调，则u是一个拟同构。

事实上，设 \(u: \mathrm{E} \to \mathrm{E}'\) 是一个复形态射。如果 Ker \(u\) （相应地 Coker \(u\) ）具有零同调，则典范态射 \(\mathrm{E} \to \operatorname{Im} u\) （相应地， \(\operatorname{Im} u \to \mathrm{E}'\) ）由推论 1 是一个拟同构。

命题 2. — 考虑一个具有正合行的复形交换图

\[
\begin{array}{c} 0 \longrightarrow \mathrm{C} ^ {\prime} \stackrel {{u}} {{\longrightarrow}} \mathrm{C} \stackrel {{v}} {{\longrightarrow}} \mathrm{C} ^ {\prime \prime} \longrightarrow 0 \\ f ^ {\prime} \Big \downarrow \qquad \qquad \qquad f \Big \downarrow \qquad \qquad f ^ {\prime \prime} \Big \downarrow \\ 0 \longrightarrow \mathrm{C} _ {1} ^ {\prime} \stackrel {{u _ {1}}} {{\longrightarrow}} \mathrm{C} _ {1} \stackrel {{v _ {1}}} {{\longrightarrow}} \mathrm{C} _ {1} ^ {\prime \prime} \longrightarrow 0. \end{array}
\]

则 \(\mathbf{H}(f')\circ \partial (u,v) = \partial (u_1,v_1)\circ \mathbf{H}(f'')\)

设 \(\alpha'' \in \mathrm{H}(\mathbf{C}'')\) ；令 \(z''\) 为类 \(\alpha''\) 的一个闭链，并令 \(x\) 为 C 的一个元素，使得 \(v(x) = z''\) ；我们有

\[
\begin{array}{r l} \big (\partial_ {u _ {1}, v _ {1}} \circ \mathrm{H} (f ^ {\prime \prime}) \big) (\alpha^ {\prime \prime}) = & \partial_ {u _ {1}, v _ {1}} (\overline {{f ^ {\prime \prime} (z ^ {\prime \prime})}}) = \overline {{u _ {1} ^ {- 1} (d f (x))}} = \overline {{f ^ {\prime} (u ^ {- 1} (d x))}} = \\ & = \mathrm{H} (f ^ {\prime}) (\overline {{u ^ {- 1} (d x)}}) = (\mathrm{H} (f ^ {\prime}) \circ \partial_ {u, v}) (\alpha^ {\prime \prime}). \end{array}
\]

例。— 设 C 是一个复形。考虑典范正合列

\[
0 \longrightarrow \mathrm{Z} (\mathrm{C}) \stackrel {j} {\longrightarrow} \mathrm{C} \stackrel {\delta} {\longrightarrow} \mathrm{B} (\mathrm{C}) (- 1) \longrightarrow 0,\tag{I}
\]

并设 \(i: \mathbf{B}(\mathbf{C}) \to \mathbf{Z}(\mathbf{C})\) 为典范嵌入。则连接同态 \(\partial(j, \delta): \mathrm{H}(\mathrm{B}(\mathrm{C})) (-1) \to \mathrm{H}(\mathrm{Z}(\mathrm{C})) (-1)\) 与 \(\mathrm{H}(i) (-1)\) 等同，这一点可以立即验证。由于 \(\mathrm{H}(\delta) = 0\) ，与（I）相关联的同调正合列分解为短正合列

\[
0 \longrightarrow \mathrm{B} _ {n} (\mathrm{C}) \stackrel {i} {\longrightarrow} Z _ {n} (\mathrm{C}) \stackrel {\mathrm{H} (j)} {\longrightarrow} \mathrm{H} _ {n} (\mathrm{C}) \longrightarrow 0.\tag{\((\mathrm{II}_{\mathfrak{n}})\)}
\]

* 应用：

\begin{enumerate}
\def\labelenumi{\arabic{enumi})}
\tightlist
\item
  奇异同调
\end{enumerate}

设 A 为一个环。对于每个拓扑空间 X，定义 X 的以 A 为系数的奇异复形 C(X, A) 如下：

在 \(\mathbf{R}^{(N)}\) 中，设 \((e_n)\) 为典范基；标准 \(n\) -单形定义为凸包 \(\Delta_n\) （即 \(\{e_0, \ldots, e_n\}\) 的凸包）。对于 \(i \in \{0, \ldots, n\}\) ，定义仿射映射 \(\iota_i: \Delta_{n-1} \to \Delta_n\) 如下： \(\iota_i(e_k) = e_k\) （当 \(k < i\) ）， \(\iota_i(e_k) = e_{k+1}\) （当 \(k \geqslant i\) ）。用 \(\mathbf{C}_n(\mathbf{X}, \mathbf{A})\) 表示 A-模 \(\mathbf{A}^{(\Sigma_n(\mathbf{X}))}\) ，其中 \(\Sigma_n(\mathbf{X})\) 是从 \(\Delta_n\) 到 \(\mathbf{X}\) 的连续映射的集合；当 \(n < 0\) 时，设 \(\mathbf{C}_n = 0\) 。对于 \(i \in \{0, \ldots, n\}\) ，定义线性映射 \(\partial_{n,i}: \mathbf{C}_n(\mathbf{X}, \mathbf{A}) \to \mathbf{C}_{n-1}(\mathbf{X}, \mathbf{A})\) 如下： \(\partial_{n,i}(e_s) = e_{s \circ \iota_i}\) （对 \(s \in \Sigma_n(\mathbf{X})\) ），并设 \(d_n = \Sigma (-1)^t \partial_{n,i}\) 。可以验证

\[
\dots \mathrm{C} _ {n} (\mathrm{X}, \mathrm{A}) \xrightarrow {d _ {n}} \mathrm{C} _ {n - 1} (\mathrm{X}, \mathrm{A}) \rightarrow \dots
\]

是一个复形。它的同调称为X以A为系数的奇异同调，记为H(X, A)或简记为H(X)。

若Y是X的子空间，则用C(X, Y, A)表示商复形C(X, A)/C(Y, A)，并用H(X, Y, A)表示其同调。由定理I可得如下正合列：

\[
\dots \rightarrow \mathrm{H} _ {n} (\mathrm{Y}, \mathrm{A}) \rightarrow \mathrm{H} _ {n} (\mathrm{X}, \mathrm{A}) \rightarrow \mathrm{H} _ {n} (\mathrm{X}, \mathrm{Y}, \mathrm{A}) \rightarrow \mathrm{H} _ {n - 1} (\mathrm{Y}, \mathrm{A}) \rightarrow \mathrm{H} _ {n - 1} (\mathrm{X}, \mathrm{A}) \rightarrow \dots
\]

\begin{enumerate}
\def\labelenumi{\arabic{enumi})}
\setcounter{enumi}{1}
\tightlist
\item
  有限胞腔复形
\end{enumerate}

设 X 是一个豪斯多夫拓扑空间。X 的一个有限胞腔分解由闭子空间的一个递增序列 \((X_{n})_{n\in\mathbb{Z}}\) 给出，满足以下条件：(i) \(X_{n} = \varnothing\) （当 n \textless{} 0 时）；

\begin{enumerate}
\def\labelenumi{(\roman{enumi})}
\setcounter{enumi}{1}
\item
  存在一个 N ，使得 \(X_{N} = X\) （因此 \(X_{n} = X\) ，只要 n \textgreater{} N）；
\item
  对所有 \(n\) ，空间 \(\mathbf{X}_{n} - \mathbf{X}_{n-1}\) 只有有限多个连通分支，称为维数为 \(n\) 的胞腔;
\item
  对于每一个n，以及对于每一个连通分支C（属于 \(X_{n} - X_{n-1}\) ），存在一个从维数为 n 的开欧几里得球 \(\dot{B}_{n}\) （TG, VI, 第 10 页）到 C 上的同胚映射，且该映射可延拓为从闭球到 X 中的连续映射。
\end{enumerate}

可以证明这些条件蕴含： \(\mathbf{H}_n(\mathbf{X}_n, \mathbf{X}_{n-1}, \mathbf{A})\) 是一个自由 A-模 \(\Gamma_n\) ，其秩等于维数为 \(n\) 的胞腔的个数，并且 \(\mathbf{H}_i(\mathbf{X}_n, \mathbf{X}_{n-1}, \mathbf{A}) = 0\) （当 \(i \neq n\) 时）。我们有 \(\mathbf{C}(\mathbf{X}_n, \mathbf{X}_{n-1}, \mathbf{A}) = \mathbf{C}(\mathbf{X}_n, \mathbf{X}_{n-2}, \mathbf{A}) / \mathbf{C}(\mathbf{X}_{n-1}, \mathbf{X}_{n-2}, \mathbf{A})\) ，由此得到一个正合列

\[
\mathrm{H} _ {n} \left(\mathrm{X} _ {n}, \mathrm{X} _ {n - 2}\right) \longrightarrow \mathrm{H} _ {n} \left(\mathrm{X} _ {n}, \mathrm{X} _ {n - 1}\right) \xrightarrow {d _ {n}} \mathrm{H} _ {n - 1} \left(\mathrm{X} _ {n - 1}, \mathrm{X} _ {n - 2}\right) \longrightarrow \mathrm{H} _ {n - 1} \left(\mathrm{X} _ {n}, \mathrm{X} _ {n - 2}\right),
\]

其中带有连接同态 \(d_{n}:\Gamma_{n}\to\Gamma_{n-1}\) 。我们有 \(d_{n}\circ d_{n+1}=0\) ，从而可以定义一个复形 \(\Gamma:\cdots\to\Gamma_{n}\xrightarrow{d_{n}}\Gamma_{n-1}\to\cdots\) .

正合列

\[
\mathrm{H} _ {n + 1} \left(\mathrm{X} _ {p}, \mathrm{X} _ {p - 1}\right)\rightarrow \mathrm{H} _ {n} \left(\mathrm{X} _ {p - 1}\right)\rightarrow \mathrm{H} _ {n} \left(\mathrm{X} _ {p}\right)\rightarrow \mathrm{H} _ {n} \left(\mathrm{X} _ {p}, \mathrm{X} _ {p - 1}\right)\rightarrow \mathrm{H} _ {n - 1} \left(\mathrm{X} _ {p - 1}\right)
\]

通过对 \(p\) 作归纳表明： \(\mathrm{H}_{n}(\mathrm{X}_{p})=0\) （当 \(p<n\) ）， \(\mathrm{H}_{n}(\mathrm{X}_{n})=\mathrm{Ker}\left(d_{n}:\Gamma_{n}\to\Gamma_{n-1}\right)\) ，并且 \(\mathrm{H}_{n}(\mathrm{X}_{p})=\mathrm{H}_{n}(\Gamma)\) （当 \(p>n\) ）。特别地， \(\mathrm{H}_{n}(\mathrm{X})\) 与 \(\mathrm{H}_{n}(\Gamma)\) 等同 .

示例。--- 考虑球面的乘积 \(\mathbf{S}_{2} \times \mathbf{S}_{2}\) 以及复射影空间 \(\mathbf{P}_{2}(\mathbf{C})\) 。设 \(b \in \mathbf{S}_{2}\) ；通过如下取定来定义胞腔分解 \((\mathrm{Y}_{n})\) （其中 \(\mathrm{Y} = \mathbf{S}_{2} \times \mathbf{S}_{2}\) ）：

\[
\mathrm{Y} _ {0} = \mathrm{Y} _ {1} = \{(b, b) \}, \quad \mathrm{Y} _ {2} = \mathrm{Y} _ {3} = (\{b \} \times \mathbf {S} _ {2}) \cup (\mathbf {S} _ {2} \times \{b \}) \quad \text {and} \quad \mathrm{Y} _ {4} = \mathbf {S} _ {2} \times \mathbf {S} _ {2};
\]

此分解有一个 0 维胞腔、两个 2 维胞腔和一个 4 维胞腔。相应复形的微分必然为零，因此 \(\mathrm{H}_{0}(\mathrm{Y})\) , \(\mathrm{H}_{2}(\mathrm{Y})\) 与 \(\mathrm{H}_{4}(\mathrm{Y})\) 分别是秩为 1、2 和 1 的自由模，并且 \(\mathrm{H}_{n}(\mathrm{Y}) = 0\) （当 \(n \notin \{0, 2, 4\}\) 时）.

我们得到一个胞腔分解 \((\mathbf{Z}_{n})\) ，即 \(\mathbf{P}_{2}(\mathbf{C})\) 的如下分解：

\[
\mathbf {Z} _ {0} = \mathbf {Z} _ {1} = \{c \}, \quad \mathbf {Z} _ {2} = \mathbf {Z} _ {3} = \mathbf {P} _ {1} (\mathbf {C}), \quad \mathbf {Z} _ {4} = \mathbf {P} _ {2} (\mathbf {C}),
\]

其中空间 \(\mathbf{P}_{1}(\mathbf{C})\) 嵌入在 \(\mathbf{P}_{2}(\mathbf{C})\) 中（TG，VIII，第20页），而 \(c\) 是 \(\mathbf{P}_{1}(\mathbf{C})\) 的一个点；这个分解有一个0维胞腔、一个2维胞腔和一个4维胞腔。这里复形的微分同样必然为零，因此可得 \(\mathbf{H}_{n}(\mathbf{P}_{2}(\mathbf{C}))\) 与 A 同构（当 \(n \in \{0, 2, 4\}\) ），否则为0。

由于所考虑的两个空间的 2 次同调模分别是秩为 2 和 1 的自由模，这两个空间不同胚。*

